\section*{Capítulo \ref{ch:g8:balanced-equations} --- Conservación de la masa y ecuaciones ajustadas}

\begin{solution}{exo:g8:balanced-equations:1}
En una \omterm{def:g7:chemical-reactions:reaction}{reacción química}, la masa total de los productos formados es
igual a la masa total de los \omterm{def:g7:chemical-reactions:reactant}{reactivos} consumidos.
\end{solution}

\begin{solution}{exo:g8:balanced-equations:2}
El coeficiente es 3. Representa 3 \omterm{def:g7:atoms-and-molecules:atom}{átomos} de carbono y $3 \times 2 = 6$
\omterm{def:g7:atoms-and-molecules:atom}{átomos} de oxígeno.
\end{solution}

\begin{solution}{exo:g8:balanced-equations:3}
Sí: 1 C y 2 O en cada miembro.
\end{solution}

\begin{solution}{exo:g8:balanced-equations:4}
\ce{2Mg + O2 -> 2MgO}: 2 Mg y 2 O en cada miembro.
\end{solution}

\begin{solution}{exo:g8:balanced-equations:5}
Nitrógeno: 2 N a la izquierda, así que 2 \ce{NH3}; entonces hay 6 H a
la derecha, así que 3
\ce{H2}: \ce{N2 + 3H2 -> 2NH3}.
\end{solution}

\begin{solution}{exo:g8:balanced-equations:6}
Por la \omterm{def:g8:balanced-equations:conservation}{conservación de la masa}: $44 - 12 = \qty{32}{g}$ de oxígeno.
\end{solution}

\begin{solution}{exo:g8:balanced-equations:7}
El gas formado se escapa a la sala, así que la balanza ya no lo pesa.
No se destruye masa: la masa que falta está en el \omterm{def:g4:air-a-mixture-of-gases:air}{aire} de la sala.
\end{solution}

\begin{solution}{exo:g8:balanced-equations:8}
Carbono: 1 en cada miembro. Hidrógeno: 4 a la izquierda, así que 2
\ce{H2O}. Oxígeno: $2 + 2 = 4$ a la derecha, así que 2 \ce{O2}:
\ce{CH4 + 2O2 -> CO2 + 2H2O}.
\end{solution}

\begin{solution}{exo:g8:balanced-equations:9}
Convertir \ce{H2O} en \ce{H2O2} cambia la sustancia: \ce{H2O2} es el
peróxido de hidrógeno, no el agua. Solo se pueden cambiar los
coeficientes:
\ce{2H2 + O2 -> 2H2O}.
\end{solution}

\begin{solution}{exo:g8:balanced-equations:10}
$20 \times 5 = 100$ \omterm{def:g7:atoms-and-molecules:molecule}{moléculas} de oxígeno; $20 \times 4 = 80$ \omterm{def:g7:atoms-and-molecules:molecule}{moléculas}
de agua.
\end{solution}

\begin{solution}{exo:g8:balanced-equations:11}
Con 2 \ce{C2H6}: 4 C, así que 4 \ce{CO2}; 12 H, así que 6 \ce{H2O};
oxígeno a la derecha, $8 + 6 = 14$, así que 7 \ce{O2}:
\ce{2C2H6 + 7O2 -> 4CO2 + 6H2O}.
\end{solution}

\begin{solution}{exo:g8:balanced-equations:12}
Sigue marcando \qty{850.0}{g}. La caja está cerrada: los \qty{0.4}{g}
de cera que han ardido se han convertido en dióxido de carbono y vapor
de agua, que siguen dentro de la caja con el oxígeno que tomaron. No ha
entrado ni salido nada.
\end{solution}

\begin{solution}{pb:g8:balanced-equations:1}
\textbf{1.} Del oxígeno del \omterm{def:g4:air-a-mixture-of-gases:air}{aire}, que se combina con el hierro: el
óxido de hierro pesa lo que el hierro más ese oxígeno.

\textbf{2.} Todo se queda dentro del frasco: el oxígeno consumido ya
estaba en el frasco, y el óxido formado se queda ahí. La masa total no
cambia.

\textbf{3.} La \omterm{def:g8:balanced-equations:conservation}{conservación de la masa}: en una reacción, la masa total
de los productos es igual a la masa total de los \omterm{def:g7:chemical-reactions:reactant}{reactivos}.

\textbf{4.} \ce{4Fe + 3O2 -> 2Fe2O3}.

\textbf{5.} \ce{3Fe + 2O2 -> Fe3O4}.

\textbf{6.} Izquierda: 4 Fe, $3 \times 2 = 6$ O. Derecha: $2 \times 2 = 4$
Fe, $2 \times 3 = 6$ O. Ajustada.

\textbf{7.} 4 \omterm{def:g7:atoms-and-molecules:atom}{átomos} de hierro por cada 3 \omterm{def:g7:atoms-and-molecules:molecule}{moléculas} de oxígeno, así que
$400 \times
\frac{3}{4} = 300$ \omterm{def:g7:atoms-and-molecules:molecule}{moléculas} de oxígeno.

\textbf{8.} $0.28 \times \frac{3}{7} = \qty{0.12}{g}$ de oxígeno.

\textbf{9.} $0.28 + 0.12 = \qty{0.40}{g}$ de óxido de hierro.

\textbf{10.} Sí: hacen falta \qty{0.12}{g} y el frasco contenía unos
\qty{0.5}{g}.

\textbf{11.} El oxígeno consumido ha salido del \omterm{def:g4:air-a-mixture-of-gases:air}{aire} del frasco (ahora
está dentro del óxido sólido), así que el frasco contiene menos gas que
antes y entra \omterm{def:g4:air-a-mixture-of-gases:air}{aire} de la sala. La lectura de la balanza sube entonces,
en la masa del \omterm{def:g4:air-a-mixture-of-gases:air}{aire} que entra.

\textbf{12.} Se tomaron \qty{0.12}{g} de oxígeno del \omterm{def:g4:air-a-mixture-of-gases:air}{aire} del frasco.
\end{solution}
